\documentclass[11pt,a4paper]{article}
% Preâmbulo comum. Não duplicar pacotes nos módulos.
\usepackage[utf8]{inputenc}
\usepackage[T1]{fontenc}
\usepackage[brazil]{babel}
\usepackage{amsmath,amssymb,amsthm,mathtools}
\usepackage{geometry}
\geometry{a4paper,margin=2.3cm}
\usepackage{microtype}
\usepackage{enumitem}
\usepackage{graphicx}
\usepackage{xcolor}
\usepackage{booktabs}
\usepackage{hyperref}
\hypersetup{colorlinks=true,linkcolor=blue!40!black,urlcolor=blue!50!black}

\setcounter{tocdepth}{2}
\numberwithin{equation}{section}

\theoremstyle{plain}
\newtheorem{teorema}{Teorema}[section]
\newtheorem{proposicao}[teorema]{Proposição}
\newtheorem{lema}[teorema]{Lema}
\newtheorem{corolario}[teorema]{Corolário}

\theoremstyle{definition}
\newtheorem{definicao}[teorema]{Definição}
\newtheorem{exemplo}[teorema]{Exemplo}
\newtheorem{exercicio}{Exercício}[section]

\theoremstyle{remark}
\newtheorem{observacao}[teorema]{Observação}
\newtheorem{contraexemplo}[teorema]{Contraexemplo}
\newtheorem{intuicao}[teorema]{Intuição}

\newcommand{\R}{\mathbb{R}}
\newcommand{\N}{\mathbb{N}}
\newcommand{\Z}{\mathbb{Z}}
\newcommand{\Q}{\mathbb{Q}}
\DeclareRobustCommand{\eps}{\varepsilon}
\DeclareRobustCommand{\vet}[1]{\vec{#1}}
\DeclareRobustCommand{\dif}{\mathrm{d}}
\DeclareRobustCommand{\abs}[1]{\left\lvert #1 \right\rvert}
\DeclareRobustCommand{\norm}[1]{\left\lVert #1 \right\rVert}
\DeclareMathOperator{\grad}{grad}
\DeclareMathOperator{\Div}{div}
\DeclareMathOperator{\rot}{rot}

\newcommand{\base}[1]{\par\medskip\noindent\textbf{Base para relembrar.} #1\par\medskip}
\newcommand{\cuidado}[1]{\par\medskip\noindent\textbf{Quando a hipótese falha.} #1\par\medskip}

\newcommand{\figura}[2]{%
  \begin{figure}[htbp]
  \centering
  \IfFileExists{#1}{%
    \includegraphics[width=0.84\textwidth]{#1}%
  }{%
    \fbox{\parbox{0.8\textwidth}{\centering Figura ainda não gerada:\\ \texttt{#1}}}%
  }
  \caption{#2}
  \end{figure}%
}

\newenvironment{esqueleto}{%
  \par\medskip\noindent\textbf{Esqueleto da demonstração.}\par\nopagebreak
  \begin{enumerate}[leftmargin=1.4em]
}{\end{enumerate}\medskip}

\title{Gabarito: quádricas, curvas no espaço e curvatura}
\author{}
\date{}
\begin{document}
\maketitle
\tableofcontents

Cada solução repete o enunciado e faz a conta até o fim. A numeração é a da nota \texttt{05\_curvas\_espaco.tex}.

\section{Básicos}

\begin{exercicio}[Básico --- B1]
Traço \(z=0\) do elipsoide \(\dfrac{x^{2}}{4}+\dfrac{y^{2}}{9}+\dfrac{z^{2}}{16}=1\).
\end{exercicio}

\noindent\textbf{Solução.}
Substituir \(z=0\) anula o terceiro termo:
\[
\frac{x^{2}}{4}+\frac{y^{2}}{9}=1.
\]
É uma elipse no plano \(z=0\), com semieixos \(2\) e \(3\): os interceptos são \(x=\pm 2\) quando \(y=0\), e \(y=\pm 3\) quando \(x=0\). A superfície é um elipsoide, e esta elipse é um traço que a identifica como limitada e oval no plano horizontal.

\begin{exercicio}[Básico --- B2]
Traços \(z=0\) dos dois hiperboloides.
\end{exercicio}

\noindent\textbf{Solução.}
Para \(x^{2}+y^{2}-z^{2}=1\), o plano \(z=0\) dá
\[
x^{2}+y^{2}=1,
\]
a circunferência unitária. O traço não é vazio. Esta é a equação canônica do hiperboloide de uma folha com \(a=b=c=1\): a superfície corta o plano \(xy\).

Para \(z^{2}-x^{2}-y^{2}=1\), o plano \(z=0\) dá
\[
-x^{2}-y^{2}=1.
\]
O primeiro membro é \(\le 0\) e o segundo membro é \(1\). Não existe \((x,y)\) real satisfazendo a equação. O traço é vazio. Além disso,
\[
z^{2}=1+x^{2}+y^{2}\ge 1,
\]
logo \(\abs{z}\ge 1\). Há uma componente em \(z\ge 1\) e outra em \(z\le -1\): hiperboloide de duas folhas. O traço vazio em \(z=0\) é o que distingue as duas folhas da folha única, cujo traço correspondente é a circunferência.

\begin{exercicio}[Básico --- B3]
Comprimento de \(\vet{r}(t)=(2\cos t,\,2\sin t,\,t)\) em \([0,2\pi]\).
\end{exercicio}

\noindent\textbf{Solução.}
A velocidade é
\[
\vet{r}'(t)=(-2\sin t,\,2\cos t,\,1).
\]
A norma é
\begin{align*}
\norm{\vet{r}'(t)}
&=\sqrt{4\sin^{2}t+4\cos^{2}t+1}
=\sqrt{4(\sin^{2}t+\cos^{2}t)+1}
=\sqrt{4+1}
=\sqrt{5}.
\end{align*}
O integrando é constante. O comprimento é
\[
L=\int_{0}^{2\pi}\sqrt{5}\,\dif t=2\pi\sqrt{5}.
\]

\begin{exercicio}[Básico --- B4]
Velocidade, aceleração e reta tangente de \(\vet{r}(t)=(t,t^{2},t^{3})\) em \(t=1\).
\end{exercicio}

\noindent\textbf{Solução.}
\begin{align*}
\vet{r}'(t)&=(1,\,2t,\,3t^{2}),\\
\vet{r}''(t)&=(0,\,2,\,6t).
\end{align*}
Em \(t=1\),
\[
\vet{r}(1)=(1,1,1),
\qquad
\vet{r}'(1)=(1,2,3),
\qquad
\vet{r}''(1)=(0,2,6).
\]
A norma da velocidade é
\[
\norm{\vet{r}'(1)}=\sqrt{1^{2}+2^{2}+3^{2}}=\sqrt{1+4+9}=\sqrt{14}\neq 0.
\]
A reta tangente é
\[
\vet{R}(s)=(1,1,1)+s(1,2,3),\qquad s\in\R,
\]
ou, em coordenadas, \(X=1+s\), \(Y=1+2s\), \(Z=1+3s\).

\section{Intermediários}

\begin{exercicio}[Intermediário --- I1]
Curvatura do círculo de raio \(a\).
\end{exercicio}

\noindent\textbf{Solução.}
Seja \(\vet{r}(t)=(a\cos t,\,a\sin t,\,0)\), com \(a>0\). Então
\begin{align*}
\vet{r}'(t)&=(-a\sin t,\,a\cos t,\,0),\\
\vet{r}''(t)&=(-a\cos t,\,-a\sin t,\,0),\\
\norm{\vet{r}'(t)}
&=\sqrt{a^{2}\sin^{2}t+a^{2}\cos^{2}t}
=a.
\end{align*}
A velocidade não se anula. O produto vetorial é
\begin{align*}
\vet{r}'\times\vet{r}''
&=
\begin{vmatrix}
\vet{i}&\vet{j}&\vet{k}\\
-a\sin t&a\cos t&0\\
-a\cos t&-a\sin t&0
\end{vmatrix}
=(0,\,0,\,a^{2}\sin^{2}t+a^{2}\cos^{2}t)
=(0,\,0,\,a^{2}).
\end{align*}
A componente \(\vet{k}\) é \((-a\sin t)(-a\sin t)-(a\cos t)(-a\cos t)=a^{2}\sin^{2}t+a^{2}\cos^{2}t=a^{2}\). As componentes \(\vet{i}\) e \(\vet{j}\) são zero porque a terceira coordenada de \(\vet{r}'\) e de \(\vet{r}''\) é zero. A norma do produto é \(\abs{a^{2}}=a^{2}\). Portanto
\[
\kappa
=\frac{\norm{\vet{r}'\times\vet{r}''}}{\norm{\vet{r}'}^{3}}
=\frac{a^{2}}{a^{3}}
=\frac{1}{a}.
\]
O raio de curvatura é \(1/\kappa=a\), o raio do círculo.

\begin{exercicio}[Intermediário --- I2]
Curvatura de \(\vet{r}(t)=(t,t^{2},t)\).
\end{exercicio}

\noindent\textbf{Solução.}
\begin{align*}
\vet{r}'(t)&=(1,\,2t,\,1),\\
\vet{r}''(t)&=(0,\,2,\,0).
\end{align*}
A norma da velocidade é
\[
\norm{\vet{r}'(t)}
=\sqrt{1+4t^{2}+1}
=\sqrt{2+4t^{2}}
=\sqrt{2(1+2t^{2})}
=\sqrt{2}\,\sqrt{1+2t^{2}}.
\]
Ela nunca se anula. O produto vetorial é
\begin{align*}
\vet{r}'\times\vet{r}''
&=
\begin{vmatrix}
\vet{i}&\vet{j}&\vet{k}\\
1&2t&1\\
0&2&0
\end{vmatrix}.
\end{align*}
Componente \(\vet{i}\): \(2t\cdot 0-1\cdot 2=-2\).

Componente \(\vet{j}\): \(-\bigl(1\cdot 0-1\cdot 0\bigr)=0\).

Componente \(\vet{k}\): \(1\cdot 2-2t\cdot 0=2\).

Logo \(\vet{r}'\times\vet{r}''=(-2,\,0,\,2)\) e
\[
\norm{\vet{r}'\times\vet{r}''}
=\sqrt{4+0+4}
=\sqrt{8}
=2\sqrt{2}.
\]
A curvatura é
\begin{align*}
\kappa(t)
&=\frac{2\sqrt{2}}{\bigl(\sqrt{2}\,\sqrt{1+2t^{2}}\bigr)^{3}}
=\frac{2\sqrt{2}}{(\sqrt{2})^{3}\,(1+2t^{2})^{3/2}}
=\frac{2\sqrt{2}}{2\sqrt{2}\,(1+2t^{2})^{3/2}}
=\frac{1}{(1+2t^{2})^{3/2}}.
\end{align*}
Usamos \((\sqrt{2})^{3}=(\sqrt{2})^{2}\cdot\sqrt{2}=2\sqrt{2}\). Em \(t=1\),
\[
\kappa(1)=\frac{1}{(1+2)^{3/2}}=\frac{1}{3^{3/2}}=\frac{1}{3\sqrt{3}}.
\]
Racionalizando,
\[
\frac{1}{3\sqrt{3}}\cdot\frac{\sqrt{3}}{\sqrt{3}}=\frac{\sqrt{3}}{9}.
\]

\begin{exercicio}[Intermediário --- I3]
\(\norm{\dif\vet{T}/\dif s}\) na hélice.
\end{exercicio}

\noindent\textbf{Solução.}
Para \(\vet{r}(t)=(\cos t,\sin t,t)\), a nota já calculou \(\norm{\vet{r}'(t)}=\sqrt{2}\) e
\[
\vet{T}(t)
=\Biggl(-\frac{\sin t}{\sqrt{2}},\,\frac{\cos t}{\sqrt{2}},\,\frac{1}{\sqrt{2}}\Biggr).
\]
O comprimento de arco a partir de \(t=0\), no sentido crescente, é
\[
s(t)=\int_{0}^{t}\sqrt{2}\,\dif u=\sqrt{2}\,t.
\]
A função \(s\) é de classe \(C^1\), estritamente crescente, e \(t=s/\sqrt{2}\). A derivada \(\dif s/\dif t=\sqrt{2}\) nunca se anula. Pela regra da cadeia,
\[
\frac{\dif\vet{T}}{\dif s}
=\frac{\dif\vet{T}/\dif t}{\dif s/\dif t}
=\frac{\vet{T}'(t)}{\sqrt{2}}.
\]
Derivando \(\vet{T}\),
\[
\vet{T}'(t)
=\Biggl(-\frac{\cos t}{\sqrt{2}},\,-\frac{\sin t}{\sqrt{2}},\,0\Biggr),
\]
logo
\[
\frac{\dif\vet{T}}{\dif s}
=\Biggl(-\frac{\cos t}{2},\,-\frac{\sin t}{2},\,0\Biggr).
\]
A norma é
\[
\norm{\frac{\dif\vet{T}}{\dif s}}
=\sqrt{\frac{\cos^{2}t}{4}+\frac{\sin^{2}t}{4}}
=\sqrt{\frac14}
=\frac12.
\]
Este é exatamente \(\kappa\). A curvatura é a taxa com que \(\vet{T}\) vira por unidade de comprimento de arco, e na hélice essa taxa vale \(1/2\) em todo ponto.

\section{Avançados}

\begin{exercicio}[Avançado --- A1]
Curvatura nula se, e somente se, a imagem está numa reta.
\end{exercicio}

\noindent\textbf{Solução.}
Hipótese comum: \(\vet{r}\) é de classe \(C^2\) num intervalo \(I\), e \(\vet{r}'(t)\neq\vet{0}\) para todo \(t\in I\).

Primeiro, suponha que a imagem está contida numa reta. Existem um ponto \(\vet{p}\) e um vetor \(\vet{d}\) com \(\norm{\vet{d}}=1\) tais que, para cada \(t\), existe um escalar \(\lambda(t)\) com
\[
\vet{r}(t)=\vet{p}+\lambda(t)\,\vet{d}.
\]
O escalar é a função de classe \(C^2\)
\[
\lambda(t)=\bigl(\vet{r}(t)-\vet{p}\bigr)\cdot\vet{d},
\]
porque \(\vet{r}\) é de classe \(C^2\). Derivando,
\[
\vet{r}'(t)=\lambda'(t)\,\vet{d},
\qquad
\vet{r}''(t)=\lambda''(t)\,\vet{d}.
\]
Os dois vetores são paralelos, logo \(\vet{r}'\times\vet{r}''=\vet{0}\). Como \(\vet{r}'\neq\vet{0}\), a fórmula da curvatura dá \(\kappa=0\) em todo ponto.

Reciprocamente, suponha \(\kappa\equiv 0\). A fórmula \(\kappa=\norm{\vet{T}'}/\norm{\vet{r}'}\) e a hipótese \(\norm{\vet{r}'}>0\) forçam \(\vet{T}'=\vet{0}\) em \(I\). Uma função de classe \(C^1\) com derivada identicamente nula num intervalo é constante: para \(t,t_{0}\in I\),
\[
\vet{T}(t)-\vet{T}(t_{0})
=\int_{t_{0}}^{t}\vet{T}'(u)\,\dif u=\vet{0}.
\]
Seja \(\vet{T}_{0}\) esse vetor constante. Por definição, \(\vet{r}'(t)=v(t)\vet{T}_{0}\) com \(v(t)=\norm{\vet{r}'(t)}>0\). Integrando a partir de um \(t_{0}\in I\),
\begin{align*}
\vet{r}(t)
&=\vet{r}(t_{0})+\int_{t_{0}}^{t}\vet{r}'(u)\,\dif u
=\vet{r}(t_{0})+\Biggl(\int_{t_{0}}^{t} v(u)\,\dif u\Biggr)\vet{T}_{0}.
\end{align*}
Todo ponto da imagem é \(\vet{r}(t_{0})\) mais um escalar vezes o vetor fixo \(\vet{T}_{0}\). A imagem está na reta que passa por \(\vet{r}(t_{0})\) com direção \(\vet{T}_{0}\).

\begin{exercicio}[Avançado --- A2]
A cúspide \(\vet{r}(t)=(t^{2},t^{3},0)\).
\end{exercicio}

\noindent\textbf{Solução.}
(i) A velocidade é \(\vet{r}'(t)=(2t,\,3t^{2},\,0)\). Em \(t=0\), \(\vet{r}'(0)=(0,0,0)\). As fórmulas de \(\vet{T}\) e de \(\kappa\) dividem por \(\norm{\vet{r}'}\) e não estão definidas nesse ponto.

A imagem: \(x=t^{2}\), \(y=t^{3}\), \(z=0\). Então \(y^{2}=t^{6}=(t^{2})^{3}=x^{3}\), com \(z=0\) e \(x\ge 0\). É a parábola semicúbica no plano \(z=0\).

(ii) Para \(t\neq 0\), \(\vet{r}''(t)=(2,\,6t,\,0)\) e
\[
\norm{\vet{r}'(t)}
=\sqrt{4t^{2}+9t^{4}}
=\sqrt{t^{2}(4+9t^{2})}
=\abs{t}\sqrt{4+9t^{2}}.
\]
O produto vetorial é
\begin{align*}
\vet{r}'\times\vet{r}''
&=
\begin{vmatrix}
\vet{i}&\vet{j}&\vet{k}\\
2t&3t^{2}&0\\
2&6t&0
\end{vmatrix}
=(0,\,0,\,12t^{2}-6t^{2})
=(0,\,0,\,6t^{2}).
\end{align*}
A componente \(\vet{k}\) é \(2t\cdot 6t-3t^{2}\cdot 2=12t^{2}-6t^{2}=6t^{2}\). Como \(t^{2}\ge 0\), a norma é \(6t^{2}\). Portanto, para \(t\neq 0\),
\begin{align*}
\kappa(t)
&=\frac{6t^{2}}{\bigl(\abs{t}\sqrt{4+9t^{2}}\bigr)^{3}}
=\frac{6t^{2}}{\abs{t}^{3}\,(4+9t^{2})^{3/2}}
=\frac{6}{\abs{t}\,(4+9t^{2})^{3/2}}.
\end{align*}
Na última passagem usamos \(t^{2}/\abs{t}^{3}=1/\abs{t}\).

(iii) O denominador tende a \(0\) quando \(t\to 0\), porque \(\abs{t}\to 0\) e \((4+9t^{2})^{3/2}\to 4^{3/2}=8\). O numerador é a constante \(6\). Logo
\[
\lim_{t\to 0}\kappa(t)=+\infty.
\]
A curvatura explode na cúspide, e o ponto \(t=0\) permanece fora do domínio da fórmula.

\begin{exercicio}[Avançado --- A3]
O triedro é ortonormal.
\end{exercicio}

\noindent\textbf{Solução.}
Hipóteses: \(\vet{r}'\neq\vet{0}\) e \(\vet{T}'\neq\vet{0}\) no ponto considerado.

Por definição, \(\vet{T}=\vet{r}'/\norm{\vet{r}'}\), logo \(\norm{\vet{T}}=1\).

De \(\vet{T}\cdot\vet{T}=1\), derivando, \(2\vet{T}\cdot\vet{T}'=0\), logo \(\vet{T}\cdot\vet{T}'=0\). A normal principal é \(\vet{N}=\vet{T}'/\norm{\vet{T}'}\), então \(\norm{\vet{N}}=1\) e
\[
\vet{T}\cdot\vet{N}
=\frac{\vet{T}\cdot\vet{T}'}{\norm{\vet{T}'}}=0.
\]

A igualdade \(\vet{B}=\vet{T}\times\vet{N}\) é a definição de \(\vet{B}\). Falta provar que essa definição produz um vetor unitário ortogonal a \(\vet{T}\) e a \(\vet{N}\).

Primeiro, \((\vet{u}\times\vet{v})\cdot\vet{u}=0\) para quaisquer \(\vet{u},\vet{v}\). Em componentes, com \(\vet{u}=(u_{1},u_{2},u_{3})\) e \(\vet{v}=(v_{1},v_{2},v_{3})\),
\begin{align*}
(\vet{u}\times\vet{v})\cdot\vet{u}
&=(u_{2}v_{3}-u_{3}v_{2})u_{1}+(u_{3}v_{1}-u_{1}v_{3})u_{2}+(u_{1}v_{2}-u_{2}v_{1})u_{3}\\
&=u_{1}u_{2}v_{3}-u_{1}u_{3}v_{2}+u_{2}u_{3}v_{1}-u_{1}u_{2}v_{3}+u_{1}u_{3}v_{2}-u_{2}u_{3}v_{1}
=0.
\end{align*}
O mesmo cálculo com o segundo fator, ou a troca \(\vet{u}\times\vet{v}=-(\vet{v}\times\vet{u})\), dá \((\vet{u}\times\vet{v})\cdot\vet{v}=0\). Aplicando a \(\vet{u}=\vet{T}\) e \(\vet{v}=\vet{N}\),
\[
\vet{B}\cdot\vet{T}=0,
\qquad
\vet{B}\cdot\vet{N}=0.
\]

A norma sai do lema da nota, \(\norm{\vet{u}\times\vet{v}}^{2}=\norm{\vet{u}}^{2}\norm{\vet{v}}^{2}-(\vet{u}\cdot\vet{v})^{2}\):
\[
\norm{\vet{B}}^{2}
=\norm{\vet{T}\times\vet{N}}^{2}
=\norm{\vet{T}}^{2}\norm{\vet{N}}^{2}-(\vet{T}\cdot\vet{N})^{2}
=1\cdot 1-0
=1.
\]
Logo \(\norm{\vet{B}}=1\). Os três vetores são unitários e dois a dois ortogonais. Junto com a definição \(\vet{B}=\vet{T}\times\vet{N}\), formam uma base ortonormal positiva.

\end{document}
